\subsection*{The learned writing interface loses usable reading capability}
Increasing the cohort from 32 to 512 questions gives 495/512 complete
pixel-exact training answers after 16,000 updates, but no exact normal
development answers (0/16 Chinese and 0/8 English; Supplementary Note 58,
Fig.~S36). All 512 training first images are exact; none of the 24 normal
development first images matches a reference. The mismatch is therefore
already present before feedback from an earlier complete generated image,
although within-image autoregression remains a possible contributor.

Scoring fixed owning checkpoints separates an erroneous learned distribution
from a decoding-only explanation. Every supplied correct development first
image is less likely than the corresponding saved emission. At the final
state, the mean best-reference minus emitted-image log probability is
$-751.24$\,nat for Chinese and $-941.82$\,nat for English. A maximum-likelihood
decoder cannot select any of these supplied references over the existing
emission. This does not rule out other correct wordings or fonts. Relative
question conditioning nevertheless transfers: 14/16 Chinese reference events
favor their owning question and all eight paired assignment contrasts are
positive. Relative preference alone does not establish useful output.

The original source reader, using the same normal and question-erased PNGs,
reproduces all 32 ILM image-memory fingerprints. A temporary helper-only
source text head then gives 12/16 literal normal answers, zero erased answers
and 11/16 source-text answers. The two date questions receive the correct
years; the learned pixel writer's corresponding forms remain damaged.
Some literal mismatches are aliases or longer relevant wording. These are
reused development cases without independent semantic grading, and the
symbolic helper generates no ILM images. The comparison localizes substantial
lost capability to the writing interface without uniquely separating its
initial query, conditional density and image-history projection.

We consequently separate glyph formation from language continuation
(Fig.~S37). A continuous CNN encoder and decoder learn training-only semantic
anchors for 4,087 distinct canonical forms. Direct decoding is exact for
4,048 forms (foreground F1 0.999940); actual-image roundtrip is exact for
3,697 (F1 0.999634). The 80.22-second primitive run peaks at 0.728\,GiB,
excluding the large contextual body and its pretraining. Semantic alignment
is weaker than reconstruction, with mean cosine 0.77735 on the first 128
forms. Clear reconstruction neither validates unseen glyphs nor demonstrates
language understanding. The proposed continuation loop preserves the native
assistant state, predicts one complete glyph and re-encodes its actual pixels.
Fitting only its readout/stopping on 508 supported questions yields readable
new complete strings corresponding to 288 items and the year 1874
(Supplementary Note 59). Nevertheless, only 1/16 normal Chinese development
images is canonical-exact, and eleven cap; the one exact full year string
occurs in training. These are partial examples, not a semantic-accuracy estimate.

A frozen intervention replaces passage dates on two templates selected from
those outputs (Fig.~S38). Eleven of fourteen prospectively fixed new dates
have clear matching output-first transcriptions, two are canonical-exact and
all fourteen stop naturally without a target length. Three damaged/wrong
digit streams remain failures. Erased-question controls give no correct date.
Six clear outputs compose full year strings absent from supported training
answers. This supports restricted input-dependent composition through actual
image feedback, while broader unseen-answer reliability remains unresolved.
